% संपूर्ण मराठी ग्रंथातील प्रकरण 47: निस्यंदने आणि निर्णेयता
% हा संपादनयोग्य स्रोतखंड आहे. संकलनासाठी संपूर्ण स्रोतसंग्रहातील
% अक्षररूपे, पूर्वभाग, चिन्हव्याख्या आणि आकृती-स्रोत आवश्यक आहेत.
% मूळ: Open Logic Project; मजकूर आणि रूपांतर CC BY 4.0.
% यंत्रानुवाद, दुरुस्ती आणि तपासणी: OpenAI Codex —
% GPT-5.6 Sol आणि GPT-6 Sol, दोन्ही Ultra effort.
% दुरुस्ती-पुनरावलोकन: GPT-6.1 Sol, Ultra effort.
\chapter{निस्यंदने आणि निर्णेयता}\label{nml:fil:chap}










\section{प्रस्तावना}\label{nml:fil:int:sec}

तर्कशास्त्राविषयी एक महत्त्वाचा प्रश्न नेहमीच
असा असतो: ते निर्णेय आहे का? म्हणजे,
“हे सूत्र वैध आहे का?” या प्रश्नाचे
उत्तर देणारी प्रभावी पद्धत आहे का?
विधानीय तर्कशास्त्र निर्णेय आहे:
सत्यता-कोष्टक रचून एखादे
सूत्र सर्वतः-सत्य आहे का हे
आपण प्रभावीपणे तपासू शकतो;
दिलेल्या सूत्र साठी असे
सत्यता-कोष्टक सांत असते.
पण मोडल सूत्र सर्व प्रतिमानांत सत्य
आहे का हे अशाच रीतीने तपासता
येईल असे स्पष्ट नाही, कारण
प्रतिमाने अनंत संख्येने आहेत.
दिलेल्या सूत्राशी संबंधित सर्व
सांत प्रतिमानांची यादी करता येते;
कारण प्रत्यक्ष त्या सूत्रामध्ये
येणाऱ्या विधानीय चले
ना जगांचे कोणते उपसंच नेमले
आहेत, एवढेच महत्त्वाचे असते.
प्राप्यता संबंध निश्चित असेल,
तर $V(p)$ ची संभाव्य वेगवेगळी
नेमणूक म्हणजे~$W$ चे सर्व
उपसंच; $\card{W} = n$ असेल,
तर असे $2^n$ उपसंच असतात.
आपल्या सूत्र~$\varphi$ मध्ये
$m$ वेगवेगळी विधानीय चले असतील, तर निश्चित जगांचा संच आणि निश्चित
प्राप्यता संबंध यांवर या चलांची $2^{nm}$ वेगवेगळी
मूल्यांकने मिळतात; इतर चलांची मूल्यांकने या सूत्रासाठी अप्रस्तुत असतात. प्रत्येक
प्रतिमानात~$\varphi$ सर्व जगांत
सत्य आहे का हे प्रत्येक जगात
$\varphi$ चे सत्यतामूल्य काढून
तपासता येते. अर्थात सर्व
संभाव्य प्राप्यता संबंधही
तपासावे लागतात; पण~$n$
जगांवर असे संबंधही सांत
संख्येनेच असतात (नेमके
$W \times W$ चे उपसंच,
म्हणजे $2^{n^2}$ संबंध).

आपल्याला \Log{K} ऐवजी
प्रतिमानांच्या एखाद्या वर्गाने
परिभाषित तर्कशास्त्रात रस
असेल (उदाहरणार्थ,
परावर्ती संक्रमक प्रतिमानांत),
तर प्राप्यता संबंध योग्य प्रकारचा
आहे का हेही तपासता यायला हवे.
आपल्याला हव्या असलेल्या चौकटी
दिलेल्या सांत मोडल सूत्रांच्या यादीने परिभाष्य
असतील तेव्हा हे करता येते
(उदाहरणार्थ, चौकटीत \Ax{T}
आणि \Ax{4} वैध आहेत का
हे तपासून). म्हणून सर्व सांत
चौकटी क्रमाने घ्याव्यात,
प्रत्येक चौकट इच्छित वर्गात
आहे का हे तपासावे,
त्या चौकटीवरील सर्व संभाव्य
प्रतिमाने नोंदवावीत आणि
प्रत्येकात $\varphi$ सत्य आहे का
हे तपासावे. नसेल, तर थांबावे:
$\varphi$ हा संबंधित प्रतिमानांच्या
वर्गात वैध नाही.

या कल्पनेत अडचण आहे:
पाहणे केव्हा थांबवता येईल,
किंवा कधी थांबवता येईलच का,
हे आपल्याला माहीत नाही.
त्या सूत्राला सांत
प्रतिप्रतिमान असेल, तर आपली
पद्धत ते शोधेल. पण सांत
प्रतिप्रतिमान नसेल, तर
उत्तर मिळणार नाही.
ते सूत्र वैध असू शकते
(म्हणजे कोणतेच प्रतिप्रतिमान
नसते), किंवा त्याला केवळ
अनंत प्रतिप्रतिमान असू शकते,
ज्यापर्यंत आपली तपासणी
कधीच पोहोचणार नाही.
प्रतिप्रतिमान असलेल्या प्रत्येक
सूत्राला सांत प्रतिप्रतिमानही
असते हे दाखवता आले, तर
ही अडचण दूर होऊ शकते.
असे असेल, तर त्या तर्कशास्त्राला
\emph{सांत प्रतिमान गुणधर्म}
आहे असे म्हणतात.

पण तर्कशास्त्राला सांत प्रतिमान
गुणधर्म आहे हे कसे दाखवायचे?
एक मार्ग म्हणजे~$\varphi$ चे
अनंत (प्रति)प्रतिमान सांत
प्रतिमानात रूपांतरित करण्याची
पद्धत शोधणे. तसे करता आले,
तर जेव्हा एखाद्या प्रतिमानात
$\varphi$ सत्य नसते, तेव्हा
त्यातून मिळालेल्या सांत
प्रतिमानातही $\varphi$ सत्य
नसेल. हे सांत प्रतिमान
सर्व सांत प्रतिमानांच्या
यादीत येईल; आणि वैध
नसलेल्या प्रत्येक सूत्राची
अवैधता आपण अखेरीस ठरवू.
मात्र सूत्र वैध असेल, तर
आपली पद्धत थांबणार नाही.
याशिवाय या पद्धतीतून
मिळणाऱ्या सांत प्रतिमानाच्या
आकाराला कमाल मर्यादा आहे
आणि ती मर्यादा केवळ
सूत्र~$\varphi$ पासून प्रभावीपणे काढता येते
हे दाखवता आले, तर
प्रतिप्रतिमान शोधताना
किती आकारापर्यंत सांत
प्रतिमाने तपासायची हे कळेल.
तेवढ्यापर्यंत प्रतिप्रतिमान
सापडले नाही, तर कोणतेच
नाही. अशा वेळी आपली
पद्धत कोणत्याही
सूत्र~$\varphi$ साठी
“$\varphi$ वैध आहे का?” हा
प्रश्न खरोखर निर्णीत करेल.

अनंत रचनांचे सांत रचनांत
रूपांतर करण्याची अनेकदा
यशस्वी होणारी रणनीती म्हणजे
संबंधित बाबतीत सारखे
वागणारे रचनेचे घटक
“एकरूप” मानणे. प्रतिमान~$\mModel{M}$
मध्ये संबंधित बाबतीत सारखे
वागणारी अनंत जगे असतील,
तर ती \emph{सर्व} जगे अशा
जगांच्या सांत संख्येतील
(कदाचित अनंत सदस्य असलेल्या)
“वर्गां”त विभागता येतील.
दुसऱ्या शब्दांत, जगांच्या
संचाचे योग्य प्रकारे विभाजन
करायचे: प्रत्येक वर्गात
कदाचित अनंत जगे असतील,
पण वर्गांची संख्या सांत असेल.
मग नवे प्रतिमान~$\mModel{M^*}$
परिभाषित करू; त्याची जगे
म्हणजे हे वर्ग. जुन्या
प्रतिमानातील सांत संख्येतील
वर्गांमुळे नव्या प्रतिमानात
सांत संख्येने जगे असतील,
म्हणजे सांत प्रतिमान मिळेल.
जग~$w$ ज्या वर्गात आहे
त्याला $[w]$ म्हणू.
जुन्या प्रतिमानात~$\varphi$ चे
प्रतिप्रतिमान असेल, तर नवे
प्रतिमानही प्रतिप्रतिमान
हवे असल्याने
$\mSat{M}{\varphi}[w]$ असेल
तेव्हा आणि तेव्हाच
$\mSat{M^*}{\varphi}[{[w]}]$
असावे. यासाठी विभाजन
आणि~$\mModel{M^*}$ चा
प्राप्यता संबंध योग्य
प्रकारे परिभाषित करावा लागतो.

हे कसे घडेल ते पाहण्यासाठी
प्रथम प्रत्येक जग प्रत्येक
जगाला प्राप्य आहे असे
समजू. मग
$\mSat{M}{\Box \psi}[w]$
असेल तेव्हा आणि तेव्हाच
एखाद्या $v \in W$ साठी
$\mSat{M}{\Box \psi}[v]$;
$\mModel{M^*}$ साठीही
$[w]$ आणि $[v]$ वापरून
तसेच होईल. पहिली कल्पना
म्हणून $P$ हा $\varphi$ मध्ये येणाऱ्या विधानचलांचा सांत संच घेऊ.
दोन जगे~$u$ आणि~$v$
एकाच वर्गात, म्हणजे
समतुल्य, असतील तेव्हा आणि
तेव्हाच प्रतिमान~$\mModel{M}$
मधील $P$ मधील सर्व विधानीय चले
वर त्यांचे सत्यत्व एकसारखे
असेल, अशी अट घेऊ.
म्हणजे प्रत्येक $p\in P$ साठी $\mSat{M}{p}[u]$
असेल तेव्हा आणि तेव्हाच
$\mSat{M}{p}[v]$.
सर्व विधानचलांसाठी $V^*(p) = \Setabs{[w]}{w\in W
  \text{ आणि }\mSat{M}{p}[w]}$ ठेवू. हा वर्गांचा प्रतिमा-संच
सर्व चलांसाठी परिभाषित आहे; मात्र खालील सत्यतेच्या जतनाचा दावा
केवळ $P$ मधील चलांनी बनलेल्या सूत्रांसाठी आहे. आपले ध्येय असे
दाखवणे आहे की
$\mSat{M}{\varphi}[w]$
असेल तेव्हा आणि तेव्हाच
$\mSat{M^*}{\varphi}[{[w]}]$.
हे आपण विगमनानेच सिद्ध
करू. आधार प्रकरण
$\varphi \equiv p$ असेल.
प्रथम $\mSat{M}{p}[w]$
असे समजा. मग व्याख्येनुसार
$[w] \in V^*(p)$;
म्हणून $\mSat{M^*}{p}[{[w]}]$.
आता $\mSat{M^*}{p}[{[w]}]$
असे समजा. याचा अर्थ
$[w] \in V^*(p)$;
म्हणजे~$w$ शी समतुल्य
अशा एखाद्या~$v$ साठी
$\mSat{M}{p}[v]$.
पण “$w$ हे~$v$ शी
समतुल्य आहे” याचा अर्थ
“$w$ आणि~$v$ येथे
नेमकी तीच
$P$ मधील विधानीय चले
सत्य आहेत”; म्हणून
$\mSat{M}{p}[w]$.
आता विगमन-पायरी पाहू;
उदाहरणार्थ $\varphi
\ident \lnot \psi$.
मग $\mSat{M}{\lnot \psi}[w]$
असेल तेव्हा आणि तेव्हाच
$\mSat/{M}{\psi}[w]$;
असेल तेव्हा आणि तेव्हाच
$\mSat/{M^*}{\psi}[{[w]}]$
(विगमन गृहीतकानुसार);
असेल तेव्हा आणि तेव्हाच
$\mSat{M^*}{\lnot \psi}[{[w]}]$.
इतर गैर-मोडल कारकांसाठीही
तसेच. $\Box$ साठीही हे
चालते: $\mSat{M^*}{\Box
  \psi}[{[w]}]$ असे समजा.
याचा अर्थ प्रत्येक $[u]$
साठी $\mSat{M^*}{\psi}[{[u]}]$.
विगमन गृहीतकानुसार प्रत्येक
$u$ साठी $\mSat{M}{\psi}[u]$.
म्हणून $\mSat{M}{\Box \psi}[w]$.
उलट, $\mSat{M}{\Box\psi}[w]$ असेल, तर प्रत्येक $u\in W$ साठी
$\mSat{M}{\psi}[u]$. विगमन गृहीतकानुसार प्रत्येक वर्ग $[u]$ येथे
$\mSat{M^*}{\psi}[{[u]}]$; म्हणून $\mSat{M^*}{\Box\psi}[{[w]}]$.
दोन्ही प्रतिमानांतील प्राप्यता संबंध येथे वैश्विक आहे.

सर्वसाधारण प्रकरणात
$\mModel{M^*}$ साठी
प्राप्यता संबंधही परिभाषित
करावा लागतो; तेव्हा
गोष्टी अधिक गुंतागुंतीच्या
होतात. प्रतिमान~$\mModel{M^*}$
च्या प्राप्यता संबंध~$R^*$
ने वरील विगमन-सिद्धता
चालण्यासाठी आवश्यक अटी
पूर्ण केल्या, तर त्या
प्रतिमानाला
\emph{निस्यंदन} म्हणू.
मग कोणतेही निस्यंदन~$\mModel{M^*}$
येथे $[w]$ वर~$\varphi$ सत्य
असेल तेव्हा आणि तेव्हाच
$\mModel{M}$ येथे~$w$
वर~$\varphi$ सत्य असेल.
पण अशी निस्यंदने
\emph{अस्तित्वात आहेत}
हेही दाखवावे लागते;
म्हणजे आवश्यक अटी
पूर्ण करणारा~$R^*$
परिभाषित करता आला पाहिजे.
हे घडण्यासाठी~$u$ आणि~$v$
ही जगे~$\varphi$ च्या सर्व उप-सूत्रे वर
एकमत असतील तेव्हा आणि तेव्हाच समतुल्य मानू.
यात $\varphi$ मध्ये येणाऱ्या विधानीय चले वरील
एकमतही समाविष्ट आहे; भाषेतील इतर सर्व चलांवरील एकमत मागितलेले नाही.
$\varphi$ ची उप-सूत्रे
सांत संख्येनेच असल्याने
निस्यंदन सांत राहील.
निस्यंदन परिभाषित
करण्याचा एकच मार्ग नाही;
निस्यंदनाचा प्राप्यता संबंध
आवश्यक गुणधर्म
(उदाहरणार्थ परावर्तित्व,
संक्रमकता इत्यादी)
राखेल यासाठी~$R^*$
ची व्याख्या कल्पकतेने
करावी लागते.













\section{पूर्वतयारी}\label{nml:fil:pre:sec}

निस्यंदनांमुळे आपल्या मोडल तर्कशास्त्रांच्या
प्रणाली निर्णेय आहेत हे सिद्ध करता येते:
त्यांना \emph{सांत प्रतिमान गुणधर्म}
आहे हे त्याने दाखवता येते. म्हणजे,
एखाद्या प्रतिमानात सत्य (असत्य)
असलेले कोणतेही सूत्र
एखाद्या \emph{सांत} प्रतिमानातही
सत्य (असत्य) असते. निस्यंदनांची
व्याख्या उपसूत्रांखाली बंद असलेल्या
सूत्रांच्या संचांच्या
सापेक्ष केली जाते.

\begin{defn}\label{nml:fil:pre:defn:modallyclosed}
  सूत्रांचा संच~$\Gamma$
  हा \emph{उपसूत्रांखाली बंद}
  असेल, जर त्यात~$\Gamma$
  मधील प्रत्येक सूत्र
  चे प्रत्येक उपसूत्र असेल.
  शिवाय, $\Gamma$ हा
  \emph{मोडलदृष्ट्या बंद}
  असेल, जर तो उपसूत्रांखाली
  बंद असेल आणि
  $\varphi \in \Gamma$ असल्यास
  $\Box\varphi,
  \Diamond\varphi \in \Gamma$
  हेही खरे असेल.
\end{defn}

उदाहरणार्थ, सूत्र~$\varphi$
दिल्यावर त्याच्या सर्व
उप-सूत्रांचा संच
उप-सूत्रे खाली बंद
असतो. $\varphi$ च्या उप-सूत्रे
च्या संचातून एखाद्या प्रतिमानाचे
निस्यंदन परिभाषित करताना
हवा असलेला गुणधर्म त्याला
असेल: मूळ प्रतिमानात~$\varphi$
सत्य (असत्य) असेल तेव्हा
आणि तेव्हाच निस्यंदनातही
ते सत्य (असत्य) असेल.

प्रतिमान~$\mModel{M}$ चे
$\Gamma$ मधून मिळणाऱ्या
निस्यंदनातील जगांचा संच
पुढील तुल्यता संबंधाच्या
सर्व तुल्यतावर्गांचा संच
म्हणून परिभाषित होतो.

\begin{defn}
$\mModel{M} =\tuple{W, R, V}$
असू द्या आणि $\Gamma$ हा
उप-सूत्रे खाली बंद
आहे असे समजा.
$\Gamma$ मधील तीच
सूत्रे सत्य ठरवणाऱ्या
कोणत्याही दोन जगांमध्ये
संबंध~$\equiv$ लागू होईल
अशी~$W$ वर त्याची व्याख्या
करू; म्हणजे:
\[u \equiv v \quad \text{तेव्हा आणि तेव्हाच} \quad
\forall \varphi \in \Gamma : \mSat{M}{\varphi}[u] \Leftrightarrow \mSat{M}{\varphi}[v].
\]
जग~$w$ चा तुल्यतावर्ग
$[w]_\equiv$, किंवा संक्षेपाने
$[w]$, म्हणजे~$w$ शी
$\equiv$-समतुल्य असलेल्या
सर्व जगांचा संच:
\[[w] = \Setabs{v}{v \equiv w}.
\]
\end{defn}

\begin{prop}
  $\mModel{M}$ आणि $\Gamma$
  दिल्यावर वरीलप्रमाणे
  परिभाषित $\equiv$ हा
  तुल्यता संबंध आहे;
  म्हणजे तो परावर्ती,
  सममित आणि संक्रमक आहे.
\end{prop}

\begin{proof}
  $\equiv$ हा परावर्ती आहे,
  कारण~$w$ येथे सत्य असलेली
  $\Gamma$ मधील सूत्रे
स्वतः येथील सूत्रांशी
नेमकी जुळतात.
  तो सममित आहे: $u$ आणि~$v$
  येथे $\Gamma$ मधील तीच
  सूत्रे सत्य असतील,
  तर~$v$ आणि~$u$
  यांच्याबाबतीतही तसेच.
  तो संक्रमकही आहे:
  $u$ आणि~$v$ येथे
  $\Gamma$ मधील तीच
  सूत्रे सत्य असतील,
  आणि~$v$ व~$w$ येथेही
  तीच असतील, तर~$u$
  व~$w$ येथेही
  $\Gamma$ मधील तीच
  सूत्रे सत्य असतील.
\end{proof}

इतर प्रत्येक तुल्यता
संबंधाप्रमाणे $\equiv$
हा~$W$ चे \emph{विभाजन}
करतो; म्हणजे~$W$ चे
परस्पर असंयुक्त उपसंच
मिळतात, आणि त्यांचा
संयोग संपूर्ण~$W$ असतो.
प्रत्येक $w \in W$
हा त्यांपैकी एका
विभागाचा, म्हणजे~$[w]$
चा, घटक असतो,
कारण $w \equiv w$.
म्हणून $[w]$ हे विभाग
सर्व~$W$ व्यापतात.
ते परस्पर असंयुक्तही
आहेत: $u \in [w]$
आणि $u \in [v]$
असतील, तर $u \equiv w$
आणि $u \equiv v$.
सममिती व संक्रमकतेवरून
$w \equiv v$; म्हणून
$[w] = [v]$.












\section{निस्यंदने}\label{nml:fil:fil:sec}

$\Gamma$ मधून मिळणारे $\mModel{M}$ चे “एकमेव” निस्यंदन
परिभाषित करण्याऐवजी कोणते प्रतिमान~$\mModel{M^*}$ हे
$\mModel{M}$ चे निस्यंदन ठरते ते परिभाषित करू.
सर्व निस्यंदनांमध्ये जगांचा संच~$W^*$ आणि मूल्यांकन~$V^*$
समान असतात. मात्र वेगवेगळ्या निस्यंदनांचे प्राप्यता
संबंध~$R^*$ वेगळे असू शकतात. निस्यंदन ठरण्यासाठी
$R^*$ ने काही अटी पूर्ण कराव्या लागतात. आपला मुख्य
निष्कर्ष, म्हणजे $\varphi \in \Gamma$ असल्यास
$\mSat{M}{\varphi}[w]$ तेव्हा आणि तेव्हाच
$\mSat{M^*}{\varphi}[{[w]}]$, सिद्ध करण्यासाठी
आपण नेमक्या याच अटी वापरू.

\begin{defn}\label{nml:fil:fil:defn:filtration}
  $\Gamma$ उपसूत्रांखाली बंद आहे आणि
  $\mModel{M} = \tuple{W, R, V}$ असे समजा.
  $\mModel{M}$ चे $\Gamma$ मधून मिळणारे
  \emph{निस्यंदन} म्हणजे पुढील अटी पूर्ण करणारे
  कोणतेही प्रतिमान $\mModel{M^*} = \tuple{W^*,R^*,V^*}$:
  \begin{enumerate}
  \item $W^* = \Setabs{[w]}{w \in W}$;
  \item \label{nml:fil:fil:defn:filtration-R}
    कोणत्याही $u,v \in W$ साठी:
    \begin{enumerate}
    \item \label{nml:fil:fil:defn:filtration-R1}
      $Ruv$ असेल, तर $R^*[u][v]$;
    \item \label{nml:fil:fil:defn:filtration-R2}
      $R^*[u][v]$ असेल, तर कोणत्याही $\Box \varphi \in \Gamma$ साठी,
      $\mSat{M}{\Box\varphi}[u]$ असल्यास $\mSat{M}{\varphi}[v]$;
    \item \label{nml:fil:fil:defn:filtration-R3}
      $R^*[u][v]$ असेल, तर कोणत्याही $\Diamond \varphi \in \Gamma$ साठी,
      $\mSat{M}{\varphi}[v]$ असल्यास $\mSat{M}{\Diamond\varphi}[u]$.
    \end{enumerate}
  \item $V^*(p) = \Setabs{[u]}{u \in V(p)}$.
  \end{enumerate}
\end{defn}

$V^*(p)$ म्हणजे काय ते नीट पाहू: $\mModel{M}$ मध्ये
ज्या प्रत्येक जग~$w$ वर~$p$ सत्य आहे, त्याच्या
तुल्यतावर्ग~$[w]$ यांचा तो संच आहे. एका बाजूने,
$w \in V(p)$ असेल, तर
या व्याख्येवरून $[w] \in V^*(p)$. पण उलट,
$[w] \in V^*(p)$ असल्यावर $w \in V(p)$ असेलच असे नाही.
$[w] \in V^*(p)$ असेल, तर \emph{एखाद्या}~$u \in V(p)$
साठी $[w] = [u]$ एवढीच खात्री मिळते. अर्थात
$[w] = [u]$ याचा अर्थ $w \equiv u$. म्हणून
$[w] \in V^*(p)$ असल्यावर एखाद्या $u \in V(p)$
साठी $w \equiv u$ आहे, एवढाच निष्कर्ष काढता येतो.

\begin{thm}\label{nml:fil:fil:thm:filtrations}
  $\mModel{M^*}$ हे $\mModel{M}$ चे $\Gamma$ मधून
  मिळणारे निस्यंदन असेल, तर प्रत्येक $\varphi \in \Gamma$
  आणि $w \in W$ साठी $\mSat{M}{\varphi}[w]$ असेल तेव्हा
  आणि तेव्हाच $\mSat{M^*}{\varphi}[{[w]}]$ असेल.
\end{thm}

\begin{proof}
  $\Gamma$ उपसूत्रांखाली बंद आहे हे वापरून~$\varphi$ वर
  विगमन करू. $\varphi \in \Gamma$ आणि $\Gamma$ उप-सूत्रे
  खाली बंद असल्याने~$\varphi$ ची सर्व उप-सूत्रे
  देखील $\in \Gamma$. त्यामुळे प्रत्येक विगमन-पायरीत
  $\varphi$ च्या उप-सूत्रे ना विगमन गृहीतक लागू होते.
  \begin{enumerate}
    \item \indcase{\varphi}{\lfalse}{$\mSat{M}{\indfrm}[w]$
        किंवा $\mSat{M^*}{\indfrm}[{[w]}]$ यांपैकी काहीच खरे नाही.}
    \item \indcase{\varphi}{\ltrue}{$\mSat{M}{\indfrm}[w]$
        आणि $\mSat{M^*}{\indfrm}[{[w]}]$ दोन्ही खरे आहेत.}
    \item \indcase{\varphi}{p}{डावीकडून उजवीकडे जाणारा निष्कर्ष
      तात्काळ मिळतो: $\mSat{M}{\indfrm}[w]$ असेल, तर
      $w \in V(p)$; त्यामुळे $[w] \in V^*(p)$,
      म्हणजेच $\mSat{M^*}{\indfrm}[{[w]}]$.
      उलट, $\mSat{M^*}{\indfrm}[{[w]}]$ असे समजा;
      म्हणजे $[w] \in V^*(p)$. मग एखाद्या
      $v \in V(p)$ साठी $w \equiv v$. म्हणून
      $\mSat{M}{p}[v]$ सुद्धा. $w \equiv v$ असल्याने
      $w$ व $v$ येथे $\Gamma$ मधील नेमकी तीच
      सूत्रे सत्य आहेत. गृहीतकानुसार
      $p \in \Gamma$ आणि $\mSat{M}{p}[v]$;
      म्हणून $\mSat{M}{\indfrm}[w]$.}

    \item \indcase{\varphi}{\lnot
          \psi}{$\mSat{M}{\indfrm}[w]$ तेव्हा आणि तेव्हाच
          $\mSat/{M}{\psi}[w]$. विगमन गृहीतकानुसार,
          $\mSat/{M}{\psi}[w]$ तेव्हा आणि तेव्हाच
          $\mSat/{M^*}{\psi}[{[w]}]$. अखेरीस,
          $\mSat/{M^*}{\psi}[{[w]}]$ तेव्हा आणि तेव्हाच
          $\mSat{M^*}{\indfrm}[{[w]}]$.}

    \item \indcase{\varphi}{(\psi \land
          \chi)}{$\mSat{M}{\indfrm}[w]$ तेव्हा आणि तेव्हाच
          $\mSat{M}{\psi}[w]$ आणि $\mSat{M}{\chi}[w]$.
          विगमन गृहीतकानुसार, $\mSat{M}{\psi}[w]$ तेव्हा आणि तेव्हाच
          $\mSat{M^*}{\psi}[{[w]}]$, तसेच $\mSat{M}{\chi}[w]$
          तेव्हा आणि तेव्हाच $\mSat{M^*}{\chi}[{[w]}]$.
          आणि $\mSat{M^*}{\indfrm}[{[w]}]$ तेव्हा आणि तेव्हाच
          $\mSat{M^*}{\psi}[{[w]}]$ आणि
          $\mSat{M^*}{\chi}[{[w]}]$.}

    \item \indcase{\varphi}{(\psi \lor
          \chi)}{$\mSat{M}{\indfrm}[w]$ तेव्हा आणि तेव्हाच
          $\mSat{M}{\psi}[w]$ किंवा $\mSat{M}{\chi}[w]$.
          विगमन गृहीतकानुसार, $\mSat{M}{\psi}[w]$ तेव्हा आणि तेव्हाच
          $\mSat{M^*}{\psi}[{[w]}]$, तसेच $\mSat{M}{\chi}[w]$
          तेव्हा आणि तेव्हाच $\mSat{M^*}{\chi}[{[w]}]$.
          आणि $\mSat{M^*}{\indfrm}[{[w]}]$ तेव्हा आणि तेव्हाच
          $\mSat{M^*}{\psi}[{[w]}]$ किंवा
          $\mSat{M^*}{\chi}[{[w]}]$.}

    \item \indcase{\varphi}{(\psi \lif
          \chi)}{$\mSat{M}{\indfrm}[w]$ तेव्हा आणि तेव्हाच
          $\mSat/{M}{\psi}[w]$ किंवा $\mSat{M}{\chi}[w]$.
          विगमन गृहीतकानुसार, $\mSat{M}{\psi}[w]$ तेव्हा आणि तेव्हाच
          $\mSat{M^*}{\psi}[{[w]}]$, तसेच $\mSat{M}{\chi}[w]$
          तेव्हा आणि तेव्हाच $\mSat{M^*}{\chi}[{[w]}]$.
          आणि $\mSat{M^*}{\indfrm}[{[w]}]$ तेव्हा आणि तेव्हाच
          $\mSat/{M^*}{\psi}[{[w]}]$ किंवा
          $\mSat{M^*}{\chi}[{[w]}]$.}

    \item \indcase{\varphi}{(\psi \liff
          \chi)}{$\mSat{M}{\indfrm}[w]$ तेव्हा आणि तेव्हाच
          $\mSat{M}{\psi}[w]$ आणि $\mSat{M}{\chi}[w]$,
          किंवा $\mSat/{M}{\psi}[w]$ आणि $\mSat/{M}{\chi}[w]$.
          विगमन गृहीतकानुसार, $\mSat{M}{\psi}[w]$ तेव्हा आणि तेव्हाच
          $\mSat{M^*}{\psi}[{[w]}]$, तसेच $\mSat{M}{\chi}[w]$
          तेव्हा आणि तेव्हाच $\mSat{M^*}{\chi}[{[w]}]$.
          आणि $\mSat{M^*}{\indfrm}[{[w]}]$ तेव्हा आणि तेव्हाच
          $\mSat{M^*}{\psi}[{[w]}]$ आणि
          $\mSat{M^*}{\chi}[{[w]}]$, किंवा
          $\mSat/{M^*}{\psi}[{[w]}]$ आणि
          $\mSat/{M^*}{\chi}[{[w]}]$.}

    \item \indcase{\varphi}{\Box
          \psi}{$\mSat{M}{\indfrm}[w]$ असे समजा.
          $\mSat{M^*}{\indfrm}[{[w]}]$ दाखवण्यासाठी
          $R^*[w][v]$ असेल असे कोणतेही~$v$ घ्या.
          \ref{nml:fil:fil:defn:filtration}\ref{nml:fil:fil:defn:filtration-R2}
          वरून $\mSat{M}{\psi}[v]$; आणि विगमन गृहीतकानुसार
          $\mSat{M^*}{\psi}[{[v]}]$. $v$ कोणतेही असल्याने
          $\mSat{M^*}{\indfrm}[{[w]}]$.

          उलट, $\mSat{M^*}{\indfrm}[{[w]}]$ असे समजा
          आणि $Rwv$ असेल असे कोणतेही~$v$ घ्या.
          \ref{nml:fil:fil:defn:filtration}\ref{nml:fil:fil:defn:filtration-R1}
          वरून $R^*[w][v]$; म्हणून
          $\mSat{M^*}{\psi}[{[v]}]$. विगमन गृहीतकानुसार
          $\mSat{M}{\psi}[v]$. $v$ कोणतेही असल्याने
          $\mSat{M}{\indfrm}[w]$.}

    \item \indcase{\varphi}{\Diamond
          \psi}{$\mSat{M}{\indfrm}[w]$ असे समजा. मग एखाद्या
          $v \in W$ साठी $Rwv$ आणि $\mSat{M}{\psi}[v]$.
          विगमन गृहीतकानुसार $\mSat{M^*}{\psi}[{[v]}]$;
          तसेच \ref{nml:fil:fil:defn:filtration}\ref{nml:fil:fil:defn:filtration-R1}
          वरून $R^*[w][v]$. म्हणून
          $\mSat{M^*}{\indfrm}[{[w]}]$.

          आता $\mSat{M^*}{\indfrm}[{[w]}]$ असे समजा.
          मग $R^*[w][v]$ आणि $\mSat{M^*}{\psi}[{[v]}]$
          असेल असा एखादा $[v] \in W^*$ आहे.
          विगमन गृहीतकानुसार $\mSat{M}{\psi}[v]$.
          \ref{nml:fil:fil:defn:filtration}\ref{nml:fil:fil:defn:filtration-R3}
          वरून $\mSat{M}{\indfrm}[w]$.}
  \end{enumerate}
\end{proof}

\begin{prob}
  \ref{nml:fil:fil:thm:filtrations} ची सिद्धता पूर्ण करा.
\end{prob}

प्रतिमानातील प्रत्येक जगावरील सत्यतेविषयीचा निष्कर्ष
प्रतिमानातील सत्यतेसाठी आणि प्रतिमानांच्या वर्गातील
वैधतेसाठीही लागू होतो.

\begin{cor}
  $\Gamma$ उपसूत्रांखाली बंद आहे असे समजा. मग:
  \begin{enumerate}
  \item $\mModel{M^*}$ हे $\mModel{M}$ चे $\Gamma$
    मधून मिळणारे निस्यंदन असेल, तर कोणत्याही
    $\varphi \in \Gamma$ साठी $\mSat{M}{\varphi}$ असेल तेव्हा
    आणि तेव्हाच $\mSat{M^*}{\varphi}$ असेल.
  \item $\mClass{C}$ हा प्रतिमानांचा वर्ग असेल आणि
    $\Gamma(\mClass{C})$ हा $\mClass{C}$ मधील प्रतिमानांच्या
    $\Gamma$-निस्यंदनांचा वर्ग असेल, तर कोणतेही
    सूत्र~$\varphi \in \Gamma$ हे $\mClass{C}$ मध्ये
    वैध असेल तेव्हा आणि तेव्हाच ते
    $\Gamma(\mClass{C})$ मध्ये वैध असेल.
  \end{enumerate}
\end{cor}














\section{निस्यंदनांची उदाहरणे}\label{nml:fil:exf:sec}

निस्यंदने अस्तित्वात आहेत हे आपण अजून दाखवलेले नाही.
पण प्रत्यक्षात कोणतेही प्रतिमान~$\mModel{M}$ घेतले,
तरी $\mModel{M}$ ची $\Gamma$ मधून मिळणारी अनेक निस्यंदने असतात.
त्यांपैकी दोन विशेष निस्यंदने पाहू: सर्वांत सूक्ष्म
आणि सर्वांत स्थूल. त्याच प्रतिमानाच्या निस्यंदनांचे
प्राप्यता संबंध वेगळे असू शकतात
(\ref{nml:fil:fil:defn:filtration} मध्ये $W^*$ आणि~$V^*$
थेट निश्चित केले आहेत). सर्वांत सूक्ष्म निस्यंदनात
संबंधित जगांच्या जोड्या शक्य तितक्या कमी असतात,
तर सर्वांत स्थूल निस्यंदनात त्या शक्य तितक्या अधिक असतात.

\begin{defn}
  $\Gamma$ उपसूत्रांखाली बंद असेल, तर प्रतिमान~$\mModel{M}$
  चे \emph{सर्वांत सूक्ष्म} निस्यंदन~$\mModel{M^*}$ पुढीलप्रमाणे
  परिभाषित होते:
  \[  R^*[u][v] \quad \text{तेव्हा आणि तेव्हाच} \quad \exists u'\in [u] \;
  \exists v' \in [v] : Ru'v'.
  \]
\end{defn}

\begin{prop}\label{nml:fil:exf:prop:finest}
  सर्वांत सूक्ष्म निस्यंदन~$\mModel{M^*}$ हे खरोखर निस्यंदन आहे.
\end{prop}

\begin{proof}
  अशा प्रकारे परिभाषित $R^*$ हा
  \ref{nml:fil:fil:defn:filtration}\ref{nml:fil:fil:defn:filtration-R}
  मधील अटी पूर्ण करतो हे तपासायचे आहे. तिन्ही अटी क्रमाने पाहू.

  $Ruv$ असेल, तर $u \in [u]$ आणि $v \in [v]$ असल्याने
  $R^*[u][v]$ सुद्धा. म्हणून
  \ref{nml:fil:fil:defn:filtration-R1} ही अट पूर्ण होते.


        \ref{nml:fil:fil:defn:filtration-R2} साठी $\Box\varphi \in \Gamma$,
        $R^*[u][v]$ आणि $\mSat{M}{\Box\varphi}[u]$ असे समजा.
        $R^*$ च्या व्याख्येवरून $u' \equiv u$ आणि $v' \equiv v$
        अशी जगे आहेत की $Ru'v'$. $u$ आणि $u'$ हे $\Gamma$
        वरील सत्यतेबाबत एकमत असल्याने
        $\mSat{M}{\Box\varphi}[u']$; म्हणून $\mSat{M}{\varphi}[v']$.
        $\Gamma$ उप-सूत्रे खाली बंद असल्याने
        $v$ आणि $v'$ यांच्यात~$\varphi$ च्या सत्यतेबाबतही
        एकमत आहे. म्हणून अपेक्षेप्रमाणे $\mSat{M}{\varphi}[v]$.


        \ref{nml:fil:fil:defn:filtration-R3} तपासण्यासाठी
        $\Diamond\varphi \in \Gamma$, $R^*[u][v]$ आणि
        $\mSat{M}{\varphi}[v]$ असे समजा. $R^*$ च्या व्याख्येवरून
        $u' \equiv u$ आणि $v' \equiv v$ अशी जगे आहेत की
        $Ru'v'$. $v$ आणि $v'$ यांच्यात~$\Gamma$ वरील
        सत्यतेबाबत एकमत आहे आणि $\Gamma$ उप-सूत्रे
        खाली बंद आहे; म्हणून $\mSat{M}{\varphi}[v']$ सुद्धा.
        त्यामुळे $\mSat{M}{\Diamond \varphi}[u']$.
        $u$ आणि $u'$ यांच्यातही $\Gamma$ वरील
        सत्यतेबाबत एकमत असल्याने
        $\mSat{M}{\Diamond \varphi}[u]$.
\end{proof}



\begin{defn}
  $\Gamma$ उपसूत्रांखाली बंद असेल, तर प्रतिमान~$\mModel{M}$
  चे \emph{सर्वांत स्थूल} निस्यंदन~$\mModel{M^*}$
  पुढीलप्रमाणे परिभाषित करतो:
  $R^*[u][v]$ तेव्हा आणि तेव्हाच

    पुढील \emph{दोन्ही} अटी पूर्ण होतात:
  \begin{enumerate}
  \item \label{nml:fil:exf:defn:coarsest-Box}$\Box\varphi \in \Gamma$
    आणि $\mSat{M}{\Box\varphi}[u]$ असेल, तर
    $\mSat{M}{\varphi}[v]$;
  \item \label{nml:fil:exf:defn:coarsest-Diamond}$\Diamond\varphi
    \in \Gamma$ आणि $\mSat{M}{\varphi}[v]$ असेल, तर
    $\mSat{M}{\Diamond\varphi}[u]$.
  \end{enumerate}
\end{defn}

\begin{prop}
  सर्वांत स्थूल निस्यंदन~$\mModel{M^*}$ हे खरोखर निस्यंदन आहे.
\end{prop}

\begin{proof}
  $R^*$ च्या व्याख्येनुसार आता फक्त $Ruv$ वरून
  $R^*[u][v]$ मिळते का हे तपासायचे आहे. म्हणून
  $Ruv$ असे समजा. $\Box\varphi \in \Gamma$
    आणि $\mSat{M}{\Box\varphi}[u]$ असे समजा; मग स्पष्टपणे
    $\mSat{M}{\varphi}[v]$, आणि
      \ref{nml:fil:exf:defn:coarsest-Box} ही अट पूर्ण होते.
    $\Diamond\varphi \in \Gamma$ आणि $\mSat{M}{\varphi}[v]$
    असे समजा. $Ruv$ असल्याने $\mSat{M}{\Diamond\varphi}[u]$
    , आणि \ref{nml:fil:exf:defn:coarsest-Diamond}
      ही अट पूर्ण होते.
\end{proof}

\begin{ex}
  $W = \PosInt$, $Rnm$ तेव्हा आणि तेव्हाच $m = n + 1$,
  आणि $V(p) = \Setabs{2n}{n \in
    \PosInt}$ असे समजा. प्रतिमान~$\mModel{M} = \tuple{W, R, V}$
  हे \ref{nml:fil:exf:fig:ex-filtration} मध्ये दाखवले आहे.
  त्याची जगे $1$, $2$ इत्यादी आहेत. प्रत्येक जगाला
  नेमके एकच जग, म्हणजे त्याचे पुढचे जग, प्राप्य आहे;
  आणि $p$ तेव्हा आणि तेव्हाच सत्य आहे जेव्हा संख्या सम असते.

  \begin{figure}
    \centering
    \begin{tikzpicture}[modal]
      \node[world] (1) [label=below:\mFalse{p}] {$1$};
      \node[world] (2) [label=below:\mTrue{p}, right=of 1] {$2$};
      \node[world] (3) [label=below:\mFalse{p}, right=of 2] {$3$};
      \node[world] (4) [label=below:\mTrue{p}, right=of 3] {$4$};
      \node[phantom] (5) [right=of 4] {};
      \draw[->] (1) to (2);
      \draw[->] (2) to (3);
      \draw[->] (3) to (4);
      \draw[dotted] (4) to (5);
    \end{tikzpicture}

    \begin{tikzpicture}[modal]
      \node[world] (1) [label=below:\mFalse{p}] {$[1]$};
      \node[world] (2) [label=below:\mTrue{p}, right=of 1] {$[2]$};
      \draw[->,bend left] (1) to (2);
      \draw[->,bend left] (2) to (1);

      \node[world] (3) [label=below:\mFalse{p},right=of 2] {$[1]$};
      \node[world] (4) [label=below:\mTrue{p}, right=of 3] {$[2]$};
      \draw[->,bend left] (3) to (4);
      \draw[->,bend left] (4) to (3);
      \draw[->,reflexive above] (4) to (4);
    \end{tikzpicture}
    \caption{एक अनंत प्रतिमान आणि त्याची निस्यंदने.}
    \label{nml:fil:exf:fig:ex-filtration}
  \end{figure}

  आता $\Gamma$ हा~$\Box p \lif p$ च्या सर्व
  उप-सूत्रांचा संच, म्हणजे
  $\{p, \Box p, \Box p \lif p\}$, असू द्या.
  $p$ सम संख्यांवरच सत्य आहे; $\Box p$ विषम
  संख्यांवरच सत्य आहे. म्हणून $\Box p \lif p$
  हे सम संख्यांवरच सत्य आहे. दुसऱ्या शब्दांत,
  प्रत्येक विषम संख्येवर $\Box p$ सत्य, पण $p$
  आणि $\Box p \lif p$ असत्य आहेत; प्रत्येक सम
  संख्येवर $p$ आणि $\Box p \lif p$ सत्य, पण
  $\Box p$ असत्य आहे. म्हणून
  $W^* = \{ [1], [2] \}$, जेथे
  $[1] = \{1, 3, 5, \dots\}$ आणि
  $[2] = \{2, 4, 6, \dots\}$.
  $2 \in V(p)$ असल्याने $[2] \in V^*(p)$;
  $1 \notin V(p)$ असल्याने $[1] \notin V^*(p)$.
  म्हणून $V^*(p) = \{[2]\}$.

  $W^*$ वर आधारलेल्या कोणत्याही निस्यंदनाच्या
  प्राप्यता संबंधात $\tuple{[1], [2]}, \tuple{[2],[1]}$
  या जोड्या असायलाच हव्यात: $R12$ असल्याने
  \ref{nml:fil:fil:defn:filtration}\ref{nml:fil:fil:defn:filtration-R1}
  वरून $R^*[1][2]$; तसेच $R23$ असल्याने
  $R^*[2][3]$, आणि $[3]=[1]$.
  त्यात $\tuple{[1],[1]}$ मात्र असू शकत नाही:
  ती जोडी असेल, तर $R^*[1][1]$ होईल;
  पण $\mSat{M}{\Box p}[1]$ आणि
  $\mSat/{M}{p}[1]$ असल्याने
  \ref{nml:fil:fil:defn:filtration-R2} शी विरोध येईल.
  $R^*[2][2]$ असावे किंवा नसावे, अशी कोणतीही
  अट नाही. म्हणून~$\mModel{M}$ ची दोन संभाव्य
  निस्यंदने आहेत; त्यांचे प्राप्यता संबंध अनुक्रमे
  \[  \{\tuple{[1],[2]}, \tuple{[2],[1]}\} \text{आणि}
  \{\tuple{[1],[2]}, \tuple{[2],[1]}, \tuple{[2],[2]}\}.
  \]
  असे आहेत. दोन्ही प्रकरणांत~$[1]$ येथे $p$ आणि
  $\Box p \lif p$ असत्य, तर $\Box p$ सत्य आहे;
  $[2]$ येथे $p$ आणि $\Box p \lif p$ सत्य,
  तर $\Box p$ असत्य आहे.
\end{ex}


\begin{prob}
  पुढील प्रतिमान~$\mModel{M} = \tuple{W, R, V}$ पाहा.
  येथे $W = \Setabs{0\sigma}{\sigma \in \Bin^*}$,
  म्हणजे~$0$ ने सुरू होणाऱ्या $0$ व~$1$ यांच्या
  सांत क्रमांचा संच; $R\sigma\sigma'$ तेव्हा आणि तेव्हाच
  $\sigma' = \sigma 0$ किंवा $\sigma' = \sigma 1$;
  तसेच $V(p) = \{0\} \cup \Setabs{\sigma
    0}{\sigma \in W}$ आणि $V(q) = \Setabs{\sigma 1}{\sigma \in
    W}$. त्याचे चित्र असे आहे:
  \begin{center}
    \begin{tikzpicture}[modal,
        node distance=2cm
      ]

      \node[world] (0) [label={below:\mTrue{p}\\\mFalse{q}},font=\tiny] {$0$};
      \node[world] (00) [label={below:\mTrue{p}\\\mFalse{q}},
        above right=of 0,font=\tiny] {$00$};
      \node[world] (000) [label={below:\mTrue{p}\\\mFalse{q}},
        right=of 00, yshift=.8cm, font=\tiny] {$000$};
      \node[phantom] (0000) [right=of 000, yshift=.5cm] {};
      \node[phantom] (0001) [right=of 000, yshift=-.5cm] {};
      \node[world] (001) [label={below:\mFalse{p}\\\mTrue{q}},
        right=of 00, yshift=-.8cm, font=\tiny] {$001$};
      \node[phantom](0010) [right=of 001, yshift=.5cm] {};
      \node[phantom](0011) [right=of 001, yshift=-.5cm] {};
      \node[world] (01) [label={below:\mFalse{p}\\\mTrue{q}},
        below right=of 0,font=\tiny] {$01$};
      \node[world] (010) [label={below:\mTrue{p}\\\mFalse{q}},
        right=of 01, yshift=.8cm,font=\tiny] {$010$};
      \node[phantom] (0100) [right=of 010, yshift=.5cm] {};
      \node[phantom] (0101) [right=of 010, yshift=-.5cm] {};
      \node[world] (011) [label={below:\mFalse{p}\\\mTrue{q}},
        right=of 01, yshift=-.8cm,font=\tiny] {$011$};
      \node[phantom] (0110) [right=of 011, yshift=.5cm] {};
      \node[phantom] (0111) [right=of 011, yshift=-.5cm] {};

      \draw[->] (0) to (00);
      \draw[->] (0) to (01);
      \draw[->] (00) to (000);
      \draw[dotted] (000) to (0000);
      \draw[dotted] (000) to (0001);
      \draw[->] (00) to (001);
      \draw[dotted] (001) to (0010);
      \draw[dotted] (001) to (0011);
      \draw[->] (01) to (010);
      \draw[dotted] (010) to (0100);
      \draw[dotted] (010) to (0101);
      \draw[->] (01) to (011);
      \draw[dotted] (011) to (0110);
      \draw[dotted] (011) to (0111);
    \end{tikzpicture}
  \end{center}
  प्रत्येक~$w$ साठी
  $\mSat/{M}{\Box(p \lor q) \lif (\Box p \lor \Box q)}[w]$
  असे आहे.

  $\Gamma$ हा~$\Box(p \lor q) \lif
  (\Box p \lor \Box q)$ च्या उप-सूत्रांचा संच
  असू द्या. $W^*$ आणि~$V^*$ काय असतील?
  $\mModel{M}$ च्या सर्वांत सूक्ष्म निस्यंदनाचा
  प्राप्यता संबंध कोणता असेल? सर्वांत स्थूल निस्यंदनाचा कोणता?
\end{prob}












\section{निस्यंदने सांत असतात}\label{nml:fil:fin:sec}

उप-सूत्रे खाली बंद असलेल्या कोणत्याही संच~$\Gamma$
साठी आपण निस्यंदनांची व्याख्या केली. केवळ त्या
व्याख्येवरून निस्यंदने सांत असतील याची खात्री मिळत नाही.
खरे तर~$\Gamma$ अनंत असेल (उदाहरणार्थ, सर्व
सूत्रांचा संच असेल), तर निस्यंदनही अनंत असू शकते.
मात्र~$\Gamma$ सांत असेल (उदाहरणार्थ, दिलेल्या
सूत्र~$\varphi$ च्या उप-सूत्रांचा संच असेल),
तर~$\Gamma$ मधून मिळणारे प्रत्येक निस्यंदनही सांत असते.

\begin{prop}\label{nml:fil:fin:prop:filt-are-finite}
  $\Gamma$ सांत असेल, तर प्रतिमान~$\mModel{M}$ चे
  $\Gamma$ मधून मिळणारे कोणतेही निस्यंदन~$\mModel{M^*}$
  सुद्धा सांत असते.
\end{prop}

\begin{proof}
  $W^*$ चा आकार म्हणजे तुल्यता संबंध~$\equiv$ अंतर्गत
  वेगवेगळ्या वर्ग~$[w]$ ची संख्या. अशा एकाच वर्गातील
  कोणतीही दोन जगे $u$, $v$, म्हणजे $u$ आणि $v$
  अशी जगे की $u \equiv v$, ही $\Gamma$ मधील सर्व सूत्रे~$\varphi$
  यांच्या सत्यतेबाबत एकमत असतात: $\varphi \in \Gamma$ असेल,
  तर~$\varphi$ हे $u$ आणि~$v$ या दोन्ही जगांवर सत्य असते
  किंवा दोन्हीवर असत्य असते. म्हणून प्रत्येक वर्ग~$[w]$
  हा~$\Gamma$ च्या एका उपसंचाशी जुळतो:~$[w]$ मधील
  जगांवर~$\varphi$ सत्य असेल अशा सर्व $\varphi \in \Gamma$ च्या संचाशी.
  दोन वेगवेगळे वर्ग $[u]$ आणि $[v]$ हे $\Gamma$ च्या
  एकाच उपसंचाशी जुळू शकत नाहीत. कारण~$u$ येथे सत्य
  असलेल्या सूत्रांचा संच आणि~$v$ येथे सत्य असलेल्या
  सूत्रांचा संच एकच असेल, तर~$u$ आणि~$v$ यांच्यात
  $\Gamma$ मधील सर्व सूत्रांच्या सत्यतेबाबत एकमत असते;
  म्हणजे $u \equiv v$. पण मग $[u] = [v]$.
  म्हणून $W^*$ पासून~$\Pow{\Gamma}$ मध्ये एकास-एक
  फल आहे आणि त्यामुळे $\card{W^*} \le \card{\Pow{\Gamma}}$.
  म्हणून~$\Gamma$ मध्ये $n$ वाक्ये असतील, तर
  $W^*$ ची गणनसंख्या~$2^n$ पेक्षा अधिक नाही.
\end{proof}













\section{\Log{K} आणि \Log{S5} यांना सांत प्रतिमान गुणधर्म आहे}\label{nml:fil:fmp:sec}

\begin{defn}
  मोडल तर्कशास्त्राच्या प्रणाली~$\Sigma$ ला
  \emph{सांत प्रतिमान गुणधर्म} आहे असे म्हणतात, जर
  सूत्र~$\varphi$ हे~$\Sigma$ च्या एखाद्या
  प्रतिमानातील एखाद्या जगावर सत्य असेल, तर~$\varphi$
  हे~$\Sigma$ च्या एखाद्या \emph{सांत} प्रतिमानातील
  एखाद्या जगावरही सत्य असेल. येथे प्रणालीचे प्रतिमान म्हणजे
  त्या प्रणालीतील प्रत्येक सूत्र प्रत्येक जगावर सत्य असलेले प्रतिमान.
\end{defn}

\begin{prop}\label{nml:fil:fmp:prop:K-fmp}
  \Log{K} ला सांत प्रतिमान गुणधर्म आहे.
\end{prop}

\begin{proof}
  \Log{K} हा वैध सूत्रांचा संच आहे;
  म्हणजे कोणतेही प्रतिमान हे~\Log{K} चे प्रतिमान आहे.
  \ref{nml:fil:fil:thm:filtrations} नुसार,
  $\mSat{M}{\varphi}[w]$ असेल, तर~$\varphi$ च्या
  उप-सूत्रांच्या संच~$\Gamma$ मधून मिळणाऱ्या
  $\mModel{M}$ च्या कोणत्याही निस्यंदनासाठी
  $\mSat{M^*}{\varphi}[{[w]}]$. कोणत्याही सूत्र
  ची उप-सूत्रे सांत संख्येनेच असतात; म्हणून
  $\Gamma$ सांत आहे. \ref{nml:fil:fin:prop:filt-are-finite}
  नुसार $\card{W^*} \le 2^n$, जेथे $n$ ही~$\Gamma$
  मधील सूत्रांची संख्या आहे. शिवाय \Log{K}
  प्रतिमानांवर कोणतेही निर्बंध घालत नसल्याने
  $\mModel{M^*}$ हे \Log{K}-प्रतिमान आहे.
\end{proof}

तर्कशास्त्र~\Log{L} याला निस्यंदनांद्वारे
सांत प्रतिमान गुणधर्म आहे हे दाखवण्यासाठी,
\Log{L}-प्रतिमानाचे निस्यंदन हे स्वतः
\Log{L}-प्रतिमान असणे आवश्यक आहे. यासाठी
अनेकदा बरेच काम करावे लागते आणि कोणतेही
निस्यंदन नेहमीच \Log{L}-प्रतिमान देत नाही.
मात्र वैश्विक प्रतिमानांच्या बाबतीत हे खरे आहे.

\begin{prop}\label{nml:fil:fmp:prop:univ-fin}
  $\mClass{U}$ हा वैश्विक प्रतिमानांचा वर्ग
  (\ref{nml:frd:es5:prop:S5=univ} पाहा) आणि
  $\mClass{U}_\mathrm{Fin}$ हा सर्व सांत वैश्विक
  प्रतिमानांचा वर्ग असू द्या. मग कोणतेही
  सूत्र~$\varphi$ हे $\mClass{U}$ मध्ये वैध असेल
  तेव्हा आणि तेव्हाच ते $\mClass{U}_\mathrm{Fin}$
  मध्ये वैध असेल.
\end{prop}

\begin{proof}
  सांत वैश्विक प्रतिमाने ही वैश्विक प्रतिमानेच
  असल्याने डावीकडून उजवीकडे जाणारा निष्कर्ष
  तात्काळ मिळतो. उलट दिशेसाठी, वैश्विक
  प्रतिमान~$\mModel{M}$ मधील एखाद्या जगावर~$w$
  $\varphi$ असत्य आहे असे समजा. $\Gamma$ मध्ये $\varphi$
  आणि त्याची सर्व उपसूत्रे घ्या; स्पष्टपणे $\Gamma$
  सांत आहे. $\mModel{M}$ चे एखादे निस्यंदन~$\mModel{M^*}$
  घ्या. \ref{nml:fil:fin:prop:filt-are-finite} नुसार
  $\mModel{M^*}$ सांत आहे; आणि
  \ref{nml:fil:fil:thm:filtrations} नुसार~$\mModel{M^*}$
  मधील $[w]$ येथे $\varphi$ असत्य आहे. $\mModel{M^*}$
  हेही वैश्विक आहे, एवढे आता पाहायचे आहे:
  कोणतीही $u$ आणि $v$ दिली, तर गृहीतकानुसार
  $Ruv$; आणि
  \ref{nml:fil:fil:defn:filtration}\ref{nml:fil:fil:defn:filtration-R}
  नुसार $R^*[u][v]$ सुद्धा.
\end{proof}

\begin{cor}\label{nml:fil:fmp:cor:S5fmp}
  \Log{S5} ला सांत प्रतिमान गुणधर्म आहे.
\end{cor}

\begin{proof}
  प्रथम, $\varphi$ हे \Log{S5} च्या कोणत्याही प्रतिमानातील एखाद्या
  जगावर सत्य आहे असे समजा. मग $\lnot\varphi$ हे \Log{S5} मध्ये
  निष्पन्न करता येण्याजोगे नसते; अन्यथा ते त्या प्रतिमानाच्या प्रत्येक जगावर
  सत्य असते. \ref{nml:com:fra:thm:generaldet} नुसार $\lnot\varphi$
  हे परावर्ती युक्लिडी प्रतिमानांच्या वर्गात वैध नसते. म्हणून $\varphi$
  हे त्या वर्गातील एखाद्या प्रतिमानाच्या एखाद्या जगावर सत्य असते.
  \ref{nml:frd:es5:prop:S5=univ} नुसार,
  $\varphi$ हे एखाद्या परावर्ती व युक्लिडी प्रतिमानातील
  एखाद्या जगावर सत्य असेल, तर ते एखाद्या वैश्विक
  प्रतिमानातील जगावर सत्य असते.
  \ref{nml:fil:fmp:prop:univ-fin} च्या सिद्धतेतील निस्यंदन वापरल्यावर
  ते एखाद्या सांत
  वैश्विक प्रतिमानातील जगावर सत्य असते
  (म्हणजे त्या प्रतिमानाचे~$\varphi$ च्या उप-सूत्रे
  च्या संचातून मिळणारे निस्यंदन).
  प्रत्येक वैश्विक प्रतिमान परावर्ती व युक्लिडीही असते;
  म्हणून $\varphi$ हे सांत परावर्ती युक्लिडी प्रतिमानातील
  एखाद्या जगावर सत्य असते.
\end{proof}

\begin{prob}
  सर्वत्र उत्तरवर्ती किंवा परावर्ती प्रतिमानाचे
  कोणतेही निस्यंदन अनुक्रमे सर्वत्र उत्तरवर्ती किंवा
  परावर्ती असते, हे दाखवा.
\end{prob}

\begin{prob}
  सममित (संक्रमक, युक्लिडी) प्रतिमानाचे
  सममित नसलेले (संक्रमक नसलेले, युक्लिडी नसलेले)
  निस्यंदन शोधा.
\end{prob}













\section{\Log{S5} निर्णेय आहे}\label{nml:fil:dec:sec}
दिलेल्या सांत रूपबंध-यादीने प्रभावीपणे परिभाषित मोडल तर्कशास्त्राच्या
प्रणाली \emph{निर्णेय} आहेत हे दाखवण्याचा
सांत प्रतिमान गुणधर्म हा सोपा मार्ग देतो.
म्हणजे, सूत्र त्या प्रणालीत
निष्पन्न करता येण्याजोगे आहे की नाही हे ठरवणारी
संगणनक्षम पद्धत मिळते.

\begin{thm}
  \Log{S5} निर्णेय आहे.
\end{thm}

\begin{proof}
  $\varphi$ दिले आहे असे समजा आणि~$\varphi$ मध्ये येणारी
  विधानीय चले ही $p_1$, \dots, $p_k$
  यांपैकीच आहेत असे समजा. प्रत्येक~$n\geq1$ साठी
  जगांचा निश्चित संच $W=\{1,\ldots,n\}$ आणि वैश्विक संबंध
  $R=W\times W$ घेऊ. $p_1$, \dots, $p_k$ यांच्या $2^{nk}$
  मूल्यांकनांची यादी सांत आहे; इतर सर्व विधानचलांना रिक्त संच नेमू.
  त्यामुळे प्रत्येक आकारासाठी तपासायच्या वैश्विक प्रतिमानांची सांत
  प्रभावी यादी मिळते.
  म्हणून दोन प्रक्रिया \emph{समांतर} चालवू शकतो:
  एका बाजूने सिद्धता पद्धतशीरपणे तयार करून
  \Log{S5} ची सर्व प्रमेये क्रमाने नोंदवायची;
  दुसऱ्या बाजूने $1$, $2$, \dots जगे असलेली सर्व
  वैश्विक प्रतिमाने क्रमाने नोंदवायची आणि त्यांपैकी
  एखाद्या प्रतिमानातील एखाद्या जगावर~$\varphi$ असत्य
  आहे का हे तपासायचे. अखेरीस या दोनपैकी एक
  प्रक्रिया उत्तर देईल. कारण
  \ref{nml:com:fra:thm:generaldet} आणि
  \ref{nml:fil:fmp:cor:S5fmp} नुसार, एकतर~$\varphi$
  निष्पन्न करता येण्याजोगे आहे किंवा ते एखाद्या सांत
  वैश्विक प्रतिमानात असत्य आहे.
\end{proof}

वरील सिद्धता \Log{S5} साठी चालते, कारण
वैश्विक प्रतिमानांची निस्यंदने आपोआप वैश्विक
असतात. परावर्तित्व आणि सर्वत्र उत्तरवर्तित्व
यांबाबतही तसेच आहे; इतर गुणधर्मांसाठी अधिक
काम करावे लागते.













\section{निस्यंदने आणि प्राप्यतेचे गुणधर्म}\label{nml:fil:acc:sec}

पाहिल्याप्रमाणे, वैश्विक, सर्वत्र उत्तरवर्ती आणि परावर्ती
प्रतिमानांची निस्यंदने अनुक्रमे वैश्विक, सर्वत्र उत्तरवर्ती
आणि परावर्ती असतात. पण सममित किंवा संक्रमक प्रतिमानाचे
प्रत्येक निस्यंदन अनुक्रमे सममित किंवा संक्रमक असेलच असे नाही.
तरी काही बाबतीत हे गुणधर्म टिकतील अशी निस्यंदने
परिभाषित करता येतात. त्यासाठी सर्वांत स्थूल निस्यंदनाच्या
व्याख्येसारखीच पद्धत वापरू, पण~$R^*$ च्या व्याख्येत
अतिरिक्त अटी घालू. $\Gamma$ उप-सूत्रे खाली बंद
आहे असे समजा. प्रतिमान~$\mModel{M} =\tuple{W, R, V}$
मधील जगांमध्ये $u$, $v$ लागू होणारे
\ref{nml:fil:acc:tab:Cn-filtrations} मधील संबंध~$C_i(u,v)$ पाहा.
या अटींच्या संयोगांवरून $R^*[u][v]$ परिभाषित करता येतो.
उदाहरणार्थ, $R^*[u][v]$ तेव्हा आणि तेव्हाच $C_1(u,v)$
ही अट पूर्ण होते असे ठरवल्यास नेमके सर्वांत स्थूल
निस्यंदन मिळते. $R^*[u][v]$ तेव्हा आणि तेव्हाच
$C_1(u,v)$ आणि $C_2(u, v)$ या दोन्ही अटी पूर्ण होतात
असे ठरवल्यास अधिक सूक्ष्म किंवा तेवढेच सूक्ष्म निस्यंदन मिळते. केवळ $C_1(u,v)$
पेक्षा $C_1(u,v)$ आणि $C_2(u,v)$ या दोन्ही अटी पूर्ण
करणाऱ्या जगांच्या जोड्या कमी किंवा तितक्याच असल्याने हे निस्यंदन
सर्वांत स्थूल निस्यंदनापेक्षा “किमान तितके सूक्ष्म” असते.

\begin{table}[ht]
  \centering
  \begin{tabular}{|ll|}
    \hline
    \multirow{2}{*}{$C_1(u,v)$:}
    &
      $\Box\varphi \in \Gamma$ आणि $\mSat{M}{\Box\varphi}[u]$ असेल,
      तर $\mSat{M}{\varphi}[v]$; आणि\\
    &
      $\Diamond\varphi \in \Gamma$ आणि $\mSat{M}{\varphi}[v]$ असेल,
      तर $\mSat{M}{\Diamond\varphi}[u]$; \\
    \hline
    \multirow{2}{*}{$C_2(u,v)$:}
    &
      $\Box\varphi \in \Gamma$ आणि $\mSat{M}{\Box\varphi}[v]$ असेल,
      तर $\mSat{M}{\varphi}[u]$; आणि\\
    &
      $\Diamond\varphi \in \Gamma$ आणि $\mSat{M}{\varphi}[u]$ असेल,
      तर $\mSat{M}{\Diamond\varphi}[v]$; \\
    \hline
    \multirow{2}{*}{$C_3(u,v)$:}
    &
      $\Box\varphi \in \Gamma$ आणि $\mSat{M}{\Box\varphi}[u]$ असेल,
      तर $\mSat{M}{\Box\varphi}[v]$; आणि\\
    &
      $\Diamond\varphi \in \Gamma$ आणि $\mSat{M}{\Diamond\varphi}[v]$ असेल,
      तर $\mSat{M}{\Diamond\varphi}[u]$; \\
    \hline
    \multirow{2}{*}{$C_4(u,v)$:}
    &
      $\Box\varphi \in \Gamma$ आणि $\mSat{M}{\Box\varphi}[v]$ असेल,
      तर $\mSat{M}{\Box\varphi}[u]$; आणि\\
    &
      $\Diamond\varphi \in \Gamma$ आणि $\mSat{M}{\Diamond\varphi}[u]$ असेल,
      तर $\mSat{M}{\Diamond\varphi}[v]$; \\
    \hline
    \end{tabular}
  \caption{निस्यंदने परिभाषित करण्यासाठी संभाव्य जगांवरील अटी.}
  \label{nml:fil:acc:tab:Cn-filtrations}
\end{table}

\begin{thm}\label{nml:fil:acc:thm:more-filtrations}
  $\mModel{M} =\tuple{W,R,V}$ हे प्रतिमान आणि
  $\Gamma$ हा उप-सूत्रे खाली बंद संच असू द्या.
  $W^*$ आणि $V^*$ यांची व्याख्या
  \ref{nml:fil:fil:defn:filtration} प्रमाणे करा. मग:
  \begin{enumerate}
  \item $R^*[u][v]$ तेव्हा आणि तेव्हाच $C_1(u, v) \land
    C_2(u,v)$ असे समजा. मग $R^*$ सममित आहे; आणि
    $\mModel{M}$ सममित असेल, तर $\mModel{M^*} =
    \tuple{W^*,R^*,V^*}$ हे निस्यंदन आहे.
  \item $R^*[u][v]$ तेव्हा आणि तेव्हाच $C_1(u, v)
    \land C_3(u,v)$ असे समजा. मग $R^*$ संक्रमक आहे;
    आणि $\mModel{M}$ संक्रमक असेल, तर
    $\mModel{M^*}=\tuple{W^*,R^*,V^*}$ हे निस्यंदन आहे.
  \item $R^*[u][v]$ तेव्हा आणि तेव्हाच $C_1(u, v) \land C_2(u,v)
    \land C_3(u,v) \land C_4(u,v)$ असे समजा. मग $R^*$
    सममित आणि संक्रमक आहे; आणि $\mModel{M}$ सममित
    व संक्रमक असेल, तर $\mModel{M^*}=\tuple{W^*,R^*,V^*}$
    हे निस्यंदन आहे.
  \item $R^*$ ची व्याख्या $R^*[u][v]$ तेव्हा आणि तेव्हाच $C_1(u,
    v) \land C_3(u,v) \land C_4(u,v)$ अशी करा. मग $R^*$
    संक्रमक आणि युक्लिडी आहे; आणि $\mModel{M}$ संक्रमक
    व युक्लिडी असेल, तर $\mModel{M^*}=\tuple{W^*,R^*,V^*}$
    हे निस्यंदन आहे.
  \end{enumerate}
\end{thm}

\begin{proof}
  \begin{enumerate}
    \item $R^*$ सममित आहे हे लगेच दिसते, कारण $C_1(u,v)
      \Leftrightarrow C_2(v,u)$ आणि $C_2(u,v) \Leftrightarrow
      C_1(v,u)$. त्यामुळे $\mModel{M}$ सममित असेल,
      तर $\mModel{M^*}$ हे $\Gamma$ मधून मिळणारे
      निस्यंदन आहे, एवढेच दाखवायचे आहे.
      $C_1(u,v)$ या अटीमुळे


          \ref{nml:fil:fil:defn:filtration-R2} आणि
          \ref{nml:fil:fil:defn:filtration-R3} या
          \ref{nml:fil:fil:defn:filtration} मधील अटी
      पूर्ण होते. म्हणून आता फक्त
      \ref{nml:fil:fil:defn:filtration}\ref{nml:fil:fil:defn:filtration-R1},
      म्हणजे $Ruv$ वरून $R^*[u][v]$ मिळते हे तपासायचे आहे.

      $Ruv$ असे समजा. $R^*[u][v]$ दाखवण्यासाठी
      $C_1(u,v)$ आणि $C_2(u,v)$ सिद्ध कराव्या लागतील.
      $C_1$ साठी: $\Box\varphi
        \in\Gamma$ आणि $\mSat{M}{\Box\varphi}[u]$ असेल,
        तर $Ruv$ असल्याने $\mSat{M}{\varphi}[v]$ सुद्धा.
         त्याचप्रमाणे,
        $\Diamond\varphi \in \Gamma$
        आणि $\mSat{M}{\varphi}[v]$ असेल, तर $Ruv$
        असल्याने $\mSat{M}{\Diamond\varphi}[u]$.
      $C_2$ साठी: $\Box\varphi \in \Gamma$
        आणि $\mSat{M}{\Box\varphi}[v]$ असेल, तर सममितीमुळे
        $Ruv$ वरून $Rvu$; म्हणून $\mSat{M}{\varphi}[u]$.
         त्याचप्रमाणे,
        $\Diamond\varphi \in\Gamma$
        आणि $\mSat{M}{\varphi}[u]$ असेल, तर
        $\mSat{M}{\Diamond\varphi}[v]$ (कारण सममितीमुळे $Rvu$).
    \item सराव.
    \item सराव.
    \item सराव.
  \end{enumerate}
\end{proof}

\begin{prob}
  \ref{nml:fil:acc:thm:more-filtrations} ची सिद्धता पूर्ण करा.
\end{prob}













\section{युक्लिडी प्रतिमानांची निस्यंदने}\label{nml:fil:euc:sec}

\ref{nml:fil:acc:sec} मधील पद्धत युक्लिडी, किंवा
सर्वत्र उत्तरवर्ती आणि युक्लिडी, प्रतिमानांसाठी
चालत नाही. \ref{nml:fil:euc:fig:ser-eucl} च्या वरच्या भागातील
प्रतिमान पाहा; ते युक्लिडी आणि सर्वत्र उत्तरवर्ती आहे.
$\Gamma = \{p, \Box p \}$ असू द्या. $\Gamma$ मधून
निस्यंदन घेतल्यावर $[w_1] = [w_3]$, कारण $w_1$
आणि $w_3$ हीच दोन जगे~$\Gamma$ वरील सत्यतेबाबत
एकमत आहेत. कोणत्याही निस्यंदनात~$\mModel{M}$
मधून वारशाने आलेले बाणही असतील; ते
\ref{nml:fil:euc:fig:ser-eucl2} मध्ये दाखवले आहेत. ते प्रतिमान
युक्लिडी नाही. शिवाय ते युक्लिडी करण्यासाठी
त्यात आणखी बाण घालता येत नाहीत. $[w_2]$
आणि $[w_4]$ यांच्यामध्ये दोन्ही दिशांनी बाण
घालावे लागतील; मग $[w_2]$ आणि~$[w_5]$
यांच्यामध्येही. पण~$w_2$ येथे $\Box p$ सत्य
असायला हवे, तर~$w_5$ येथे $p$ असत्य आहे.

\begin{figure}[htpb]
  \centering
  \begin{tikzpicture}[modal]
    \node[world] (w1)
          [label={left: \mFalse{p}},
            label={below: $\mSat{{}}{\Box p}$}] {$w_1$};
    \node[world] (w2)
          [label={right: \mTrue{p}},
            label = {below: $\mSat{{}}{\Box p}$},
            right=of w1] {$w_2$};
    \draw[reflexive above] (w2) to (w2);
    \draw[->] (w1) to (w2);
    \node[world] (w3)
          [label={left: \mFalse{p}},
            label={below: $\mSat{{}}{\Box p}$},
            below=of w1] {$w_3$};
    \node[world] (w4)
          [label={right:\mTrue{p}},
            label={below: $\mSat/{{}}{\Box p}$},
            right=of w3] {$w_4$};
    \draw[->] (w3) to (w4);
    \draw[reflexive above] (w4) to (w4);
    \node[world] (w5)
          [label={right: \mFalse{p}},
            label=below:$\mSat/{{}}{\Box p}$,
            right=of w4] {$w_5$};
    \draw[->, bend left] (w4) to (w5);
    \draw[->, bend left] (w5) to (w4);
    \draw[reflexive above] (w5) to (w5);
  \end{tikzpicture}
  \caption{सर्वत्र उत्तरवर्ती आणि युक्लिडी प्रतिमान.}
  \label{nml:fil:euc:fig:ser-eucl}
\end{figure}

\begin{figure}[ht]
  \centering
  \begin{tikzpicture}[modal,every text node part/.style={align=center}]
    \node[world] (w1)
          [label={left: \mFalse{p}},
            label={right:$[w_1]=[w_3]$},
            label={below: $\mSat{{}}{\Box p}$}] {$[w_1]$};
    \node[world] (w2)
         [label={right: \mTrue{p}},
           label = {below: $\mSat{{}}{\Box p}$},
           right=of w1, yshift=1.5cm] {$[w_2]$};
    \draw[reflexive above] (w2) to (w2);
    \draw[->] (w1) to (w2);
    \node[world] (w4)
         [label={right:\mTrue{p}},
           label={below: $\mSat/{{}}{\Box p}$},
           right=of w1,yshift=-1.5cm] {$[w_4]$};
    \draw[->] (w1) to (w4);
    \draw[reflexive above] (w4) to (w4);
    \node[world] (w5)
         [label={right: \mFalse{p}},
             label=below:$\mSat/{{}}{\Box p}$,
             right=of w4] {$[w_5]$};
    \draw[->, bend left] (w4) to (w5);
    \draw[->, bend left] (w5) to (w4);
    \draw[reflexive above] (w5) to (w5);
  \end{tikzpicture}
  \caption{\ref{nml:fil:euc:fig:ser-eucl} मधील प्रतिमानाचे निस्यंदन.}
  \label{nml:fil:euc:fig:ser-eucl2}
\end{figure}

विशेषतः युक्लिडी निस्यंदन मिळवण्यासाठी
उप-सूत्रे खाली बंद असलेल्या मनमानी $\Gamma$
मधून निस्यंदन घेणे पुरेसे नाही. त्याऐवजी
\emph{मोडलदृष्ट्या बंद} असलेले संच~$\Gamma$
घ्यावे लागतात (\ref{nml:fil:pre:defn:modallyclosed} पाहा).
असे रिक्त नसलेले वाक्यांचे संच अनंत असतात; म्हणून त्यांच्यापासून
संबंधित प्रणालीचा सांत प्रतिमान गुणधर्म किंवा
निर्णेयता तात्काळ मिळत नाही.

\begin{thm}\label{nml:fil:euc:thm:modal-closed-filt}
  $\Gamma$ मोडलदृष्ट्या बंद असू द्या,
  $\mModel{M}=\tuple{W,R,V}$ हे प्रतिमान असू द्या,
  आणि $\mModel{M^*} = \tuple{W^*,R^*,V^*}$ हे
  $\mModel{M}$ चे सर्वांत स्थूल निस्यंदन असू द्या.
  \begin{enumerate}
  \item $\mModel{M}$ सममित असेल, तर $\mModel{M^*}$ सुद्धा सममित आहे.
  \item $\mModel{M}$ संक्रमक असेल, तर $\mModel{M^*}$ सुद्धा संक्रमक आहे.
  \item $\mModel{M}$ युक्लिडी असेल, तर $\mModel{M^*}$ सुद्धा युक्लिडी आहे.
  \end{enumerate}
\end{thm}

\begin{proof}
\begin{enumerate}
  \item सराव. \Ax{B} आणि $\Ax{B_\Diamond}$ ही दोन्ही
    सर्व सममित प्रतिमानांत वैध आहेत, हे वापरा.
  \item $\mModel{M^*}$ हे सर्वांत स्थूल निस्यंदन असेल,
    तर व्याख्येनुसार $R^*[u][v]$ तेव्हा आणि तेव्हाच
    $C_1(u,v)$. संक्रमकतेसाठी $C_1(u,v)$ आणि
    $C_1(v,w)$ असे समजा; $C_1(u,w)$ दाखवायचे आहे.
    $\Box\varphi\in\Gamma$ आणि
      $\mSat{M}{\Box \varphi}[u]$ असे समजा.
      सर्व संक्रमक प्रतिमानांत \Ax{4} वैध असल्याने
      $\mSat{M}{\Box\Box\varphi}[u]$. मोडल बंदतेवरून
      $\Box\Box\varphi \in \Gamma$; म्हणून $C_1(u,v)$ वरून
      $\mSat{M}{\Box\varphi}[v]$, आणि $C_1(v,w)$ वरून
      $\mSat{M}{\varphi}[w]$.
      $\Diamond\varphi\in\Gamma$ आणि $\mSat{M}{\varphi}[w]$ असे समजा.
      उपसूत्रांखालील बंदतेवरून $\varphi\in\Gamma$; मोडल बंदतेवरून
      $\Diamond \varphi \in \Gamma$; म्हणून $C_1(v,w)$ वरून
      $\mSat{M}{\Diamond \varphi}[v]$. तसेच मोडल बंदतेवरून
      $\Diamond\Diamond\varphi \in \Gamma$; म्हणून $C_1(u,v)$ वरून
      $\mSat{M}{\Diamond\Diamond \varphi}[u]$.
      सर्व संक्रमक प्रतिमानांत $\Ax{4}_\Diamond$
      वैध असल्याने $\mSat{M}{\Diamond\varphi}[u]$.
  \item सराव. \Ax{5} आणि $\Ax{5_\Diamond}$ ही दोन्ही
    सर्व युक्लिडी प्रतिमानांत वैध आहेत, हे वापरा.
\end{enumerate}
\end{proof}

\begin{prob}
  \ref{nml:fil:euc:thm:modal-closed-filt} ची सिद्धता पूर्ण करा.
\end{prob}



